\section{Compatible smash products and quantitative nilpotence}
\label{sec:filtered-products}

The topological argument will produce a complex-cobordism class which
vanishes on a finite pointed subspace.  We need the corresponding stable
cohomotopy class to vanish on a union of products.  The compatibility on
the intersections of these products is essential.  The following result
provides it with a fixed positive proportion of restricted factors.

Our goal is \Cref{thm:filtered-product}.  We first build a lift through
a power of the unit fiber on the entire union.  We then use quantitative
nilpotence at finitely many primes and a separate connectivity estimate
at all larger primes.  The latter step makes the eventual product
threshold independent of the prime.

All constructions with spectra in this section are derived, and a
commutative ring spectrum means an $\mathbb E_\infty$ ring spectrum.  A spectrum
is $b$-connective if its homotopy groups below degree $b$ vanish.  We write
$[K,Y]$ for the group of homotopy classes of maps of spectra and
$\operatorname{Map}(K,Y)$ for their mapping space.  For a ring spectrum
$R$ under the sphere, put
\[
I_R=\operatorname{fib}(\mathbb S\longrightarrow R),
\qquad
\varphi_{R,n}:I_R^{\wedge n}\longrightarrow\mathbb S.
\]
The second map is the smash product of the $n$ fiber inclusions, with
$I_R^{\wedge0}=\mathbb S$.  In a prime-local category, the sphere in this
notation is $\mathbb S_{(p)}$.

We use the commutative Thom-spectrum model of $MU$ and its canonical
complex orientation as recalled in
\cite[Lecture~6, Remarks~5 and~7 and Theorem~8]{Lurie2010}.

A map $K\to Y$ has $R$-Adams filtration at least $n$ if it lifts through
\[
\varphi_{R,n}\wedge\operatorname{id}_Y:
I_R^{\wedge n}\wedge Y\longrightarrow Y.
\]

\subsection{A compatible lift on the union}

A map to the sphere which becomes null after applying the unit
$\mathbb S\to R$ lifts through $I_R$.  On a product with $n$
restricted factors, this suggests a lift through $I_R^{\wedge n}$.
The difficulty is making these lifts agree where different choices
of restricted factors meet.  We construct one lift on the whole union.

\begin{proposition}[Compatible Adams lift]
\label{prop:compatible-adams-lift}
Let $R$ be a commutative ring spectrum, let $G\subset E$ be a finite
pointed CW pair, let $d\in\mathbb Z$, and let
$v:\Sigma^\infty E\to\Sigma^d\mathbb S$ be a stable map whose
restriction to $\Sigma^\infty G$ becomes null in $\Sigma^dR$.
For integers $1\leq n\leq\ell$, let
$U_{\ell,n}\subset E^{\wedge\ell}$ be the union requiring at least
$n$ factors in $G$.  Then the restriction of $v^{\wedge\ell}$ admits
a lift
\[
\Sigma^\infty U_{\ell,n}\longrightarrow\Sigma^{d\ell}I_R^{\wedge n}
\]
through $\Sigma^{d\ell}\varphi_{R,n}$.
\end{proposition}

We describe the finite diagram in which the lift will be constructed.
Let $\mathcal P_n^\times$ be the poset of nonempty subsets of
$\{1,\ldots,n\}$.  Give a subset $S$ the value $R^{\wedge|S|}$,
with maps inserting units along subset inclusions.  Adjoining the
sphere at the empty subset gives the $n$-fold cube of units.  The fiber
of its unit-cone map
\[
\mathbb S\longrightarrow
\operatorname*{holim}_{S\in\mathcal P_n^\times}R^{\wedge|S|}
\]
is $I_R^{\wedge n}$, as we verify at the end of the proof.
It therefore suffices to null the composite of
$v^{\wedge\ell}|_{U_{\ell,n}}$ with the $d\ell$-fold suspension
of this map.

We will produce this nullhomotopy by introducing simplex parameters.
The given $R$-nullhomotopy on $G$ will supply compatible nullhomotopies
at the vertices of every simplex.  We distribute one simplex parameter
among the $\ell$ product factors so that in degree $j$ at most $j$
of the resulting parameters are away from vertices.  When $j<n$, at
least one of the $n$ restricted factors then receives a vertex and
is collapsed.  The following construction makes this distribution
natural in the simplex.  The relation between unit cubes and the
cosimplicial cobar construction is recalled in
\cite[Section~2.1, Propositions~2.11 and~2.14]{MathewNaumannNoel2017}.

Let $\Delta$ denote the simplex category, with objects
$[j]=\{0,\ldots,j\}$ and monotone maps.  The topological simplices
$\Delta^j$ form a covariant functor: a monotone map sends a distribution
on its source vertices to the pushforward distribution on its target
vertices.  Write $V^j\subset\Delta^j$ for the set of vertices.

\begin{lemma}[Slicing a simplex]
\label{lem:simplex-slicing}
For every $\ell\geq1$ there is a natural continuous map
\[
\rho=(\rho_1,\ldots,\rho_\ell):
\Delta^j\longrightarrow(\Delta^j)^\ell
\]
in the variable $[j]\in\Delta$, such that at most $j$ of the points
$\rho_r(x)$ are nonvertices, for each $x\in\Delta^j$.
\end{lemma}

\begin{proof}
For $x=(x_0,\ldots,x_j)$ put
$s_a=\sum_{b<a}x_b$, so that $s_0=0$ and $s_{j+1}=1$.
For $1\leq r\leq\ell$ define
\begin{equation}
\label{eq:simplex-slicing-formula}
\rho_r(x)_a
=\ell\max\left\{0,
 \min\left(s_{a+1},\frac r\ell\right)
 -\max\left(s_a,\frac{r-1}\ell\right)\right\},
\qquad 0\leq a\leq j.
\end{equation}
This is $\ell$ times the length of the intersection of
$[s_a,s_{a+1}]$ with the $r$th interval of length $1/\ell$.
The coordinates are continuous, nonnegative, and sum to one.

Every nonempty fiber of a monotone map $\alpha:[j]\to[k]$ is a
consecutive block.  Passing from $x$ to $\alpha_*x$ merges the adjacent
mass intervals in that block.  The intersection length of the merged
interval with a fixed slice equals the sum of its constituent
intersection lengths.  Empty fibers contribute zero.  It follows that
$\alpha_*\rho_r(x)=\rho_r(\alpha_*x)$ for every $r$.  This verifies
naturality for all maps of $\Delta$, including its surjections.

A nonvertex value $\rho_r(x)$ requires at least one of the internal
breaks $s_1,\ldots,s_j$ to lie strictly between
$(r-1)/\ell$ and $r/\ell$.  Distinct such slices have disjoint interiors
and therefore require distinct breaks.  There are at most $j$ of them.
Breaks lying at slice endpoints do not create nonvertex values.
\end{proof}

\begin{figure}[ht]
\centering
\setlength{\unitlength}{0.9mm}
\begin{picture}(150,37)
\put(0,28){\makebox(26,0)[r]{mass intervals}}
\put(32,28){\line(1,0){100}}
\multiput(32,26)(100,0){2}{\line(0,1){4}}
\put(62,26){\line(0,1){4}}
\put(102,26){\line(0,1){4}}
\put(47,32){\makebox(0,0){$x_0=3/10$}}
\put(82,32){\makebox(0,0){$x_1=2/5$}}
\put(117,32){\makebox(0,0){$x_2=3/10$}}
\multiput(62,15)(0,3){4}{\line(0,1){1.5}}
\multiput(102,15)(0,3){4}{\line(0,1){1.5}}
\put(0,13){\makebox(26,0)[r]{four slices}}
\put(32,13){\line(1,0){100}}
\multiput(32,11)(25,0){5}{\line(0,1){4}}
\put(44.5,6){\makebox(0,0){vertex}}
\put(69.5,6){\makebox(0,0){nonvertex}}
\put(94.5,6){\makebox(0,0){nonvertex}}
\put(119.5,6){\makebox(0,0){vertex}}
\end{picture}
\caption{Simplex slicing for $j=2$ and $\ell=4$.  Each slice records
the proportions of the three mass intervals that it meets.  A
nonvertex value requires an internal mass break in the interior of
that slice, so there can be at most $j$ such values.}
\label{fig:simplex-slicing}
\end{figure}

\Cref{fig:simplex-slicing} shows the counting mechanism in
\Cref{lem:simplex-slicing}.  We now combine this distribution with
one chosen nullhomotopy to construct the lift.

\begin{proof}[Proof of \Cref{prop:compatible-adams-lift}]
Work in the stable functor category
$\mathcal D=\operatorname{Fun}(\Delta,\mathrm{Sp})$.
Its equivalences, cofibers, and smash products are computed objectwise.
Consider the cosimplicial commutative ring spectrum
\[
B^j=R^{\wedge(j+1)}.
\]
For a monotone map $\alpha:[j]\to[k]$, its structure map multiplies
the factors in each fiber of $\alpha$ and inserts a unit for each
empty fiber.  This describes its cofaces and codegeneracies, together
with their coherences.  In particular, the maps
$(B^j)^{\wedge\ell}\to B^j$ given by multiplication are natural in
$[j]$.

\smallskip
\noindent\emph{The nullhomotopy on the vertices.}
The functor $V$ is the covariant representable
$\operatorname{Hom}_\Delta([0],-)$.  Thus
$\Sigma^\infty(G\wedge V^\bullet_+)$ is the left Kan extension of
$\Sigma^\infty G$ from the object $[0]$.  The resulting adjunction is
an equivalence of mapping spaces
\begin{equation}
\label{eq:vertices-yoneda}
\operatorname{Map}_{\mathcal D}
 \bigl(\Sigma^\infty(G\wedge V^\bullet_+),\Sigma^dB^\bullet\bigr)
\simeq
\operatorname{Map}_{\mathrm{Sp}}(\Sigma^\infty G,\Sigma^dR).
\end{equation}
Indeed, the left Kan extension in degree $j$ is the wedge of copies of
$\Sigma^\infty G$ indexed by the $j+1$ maps $[0]\to[j]$; the
adjunction says that a map from this diagram is determined coherently
by its value on the copy in degree zero.

There is a natural map
\[
a:\Sigma^\infty(E\wedge\Delta^\bullet_+)
 \longrightarrow\Sigma^dB^\bullet
\]
which forgets the simplex coordinate and applies $v$ followed by the
unit.  Its restriction to $G\wedge V^\bullet_+$ corresponds under
\Cref{eq:vertices-yoneda} to the assumed null map into $\Sigma^dR$.
Choose a nullhomotopy there.  The mapping-space adjunction supplies a
single natural nullhomotopy on the vertex diagram.  The cofiber
universal property therefore gives a map
\[
\psi:Q^\bullet\longrightarrow\Sigma^dB^\bullet,
\qquad
Q^j=\Sigma^\infty\left(
 \frac{E\wedge\Delta^j_+}{G\wedge V^j_+}\right),
\]
and a homotopy $\psi q\simeq a$ in $\mathcal D$, where $q$ is the
quotient map.  The displayed ordinary quotients compute the derived
cofibers: the inclusion of the indicated pointed CW subspace is a
cofibration in each degree.  This also permits us to use the actual
collapse maps in the next step.

\smallskip
\noindent\emph{Collapse on the whole union.}
Put $U=U_{\ell,n}$.  The maps of \Cref{lem:simplex-slicing} give a
natural map of pointed spaces
\[
U\wedge\Delta^j_+\longrightarrow(E\wedge\Delta^j_+)^{\wedge\ell},
\quad
(x_1\wedge\cdots\wedge x_\ell,x)
\longmapsto
\bigwedge_{r=1}^\ell(x_r,\rho_r(x)).
\]
Follow it by the smash of the quotient maps to obtain, after
suspension, a map into $(Q^j)^{\wedge\ell}$.  If $j<n$, at least $n$
of the factors $x_r$ belong to $G$, while at most $j$ of the simplex
parameters $\rho_r(x)$ are nonvertices.  Some factor is consequently
in $G\wedge V^j_+$ and is collapsed.  The resulting map
\begin{equation}
\label{eq:sliced-union-collapse}
\Sigma^\infty(U\wedge\Delta^j_+)
 \longrightarrow(Q^j)^{\wedge\ell}
\end{equation}
is the zero map for every $j<n$.  It is the zero natural
transformation on these degrees, since it is already a literal
collapse on the diagrams of spaces.

Compose \Cref{eq:sliced-union-collapse} with $\psi^{\wedge\ell}$
and multiplication to obtain a natural map into
$\Sigma^{d\ell}B^j$.  Smashing the chosen factorization homotopy
$\psi q\simeq a$ and then multiplying identifies this map with the
map which forgets the simplex coordinate and applies
$v^{\wedge\ell}|_U$ followed by the unit.  This identification is
natural: each map $a$ forgets its own simplex parameter, so it is
independent of the particular values of $\rho_r(x)$.

The natural projection
\[
\Sigma^\infty(U\wedge\Delta^\bullet_+)
 \longrightarrow\operatorname{const}(\Sigma^\infty U)
\]
is an objectwise equivalence, since every $\Delta^j$ is contractible.
Invert this equivalence in the functor category and restrict to the
full subcategory $\Delta_{\leq n-1}$.  We have proved that the cone
\[
\operatorname{const}(\Sigma^\infty U)
 \longrightarrow\operatorname{const}(\Sigma^{d\ell}\mathbb S)
 \longrightarrow\Sigma^{d\ell}B^\bullet
\]
is null on that subcategory.  Inverting an objectwise equivalence here
does not require a natural choice of a simplex vertex or a natural
contraction of the simplices.

\smallskip
\noindent\emph{Restriction to the unit cube.}
The rule
\[
S\longmapsto[|S|-1]
\]
defines a functor from the punctured cube
$\mathcal P_n^\times$ to $\Delta_{\leq n-1}$:
an inclusion of subsets induces the ordered injection on their
enumerations.  Restricting $B^\bullet$ gives the diagram
$S\mapsto R^{\wedge|S|}$, whose maps insert units.  The null cone
just constructed therefore yields a null composite
\begin{equation}
\label{eq:unit-cube-null-cone}
\Sigma^\infty U\xrightarrow{v^{\wedge\ell}}
\Sigma^{d\ell}\mathbb S
\longrightarrow
\Sigma^{d\ell}\operatorname*{holim}_{S\in\mathcal P_n^\times}
 R^{\wedge|S|}.
\end{equation}
Only restriction of cones is used; no assertion about cofinality or
an identification of partial totalizations is needed.

Adjoin the value $\mathbb S$ at the empty subset.  The resulting
$n$-cube is the external smash product of $n$ copies of
$\mathbb S\to R$.  Its total fiber is $I_R^{\wedge n}$.
To verify this, compute fibers successively in its $n$ directions.
Taking the fiber in the last direction replaces every map
$C\wedge\mathbb S\to C\wedge R$ by $C\wedge I_R$, because
smash with $C$ preserves finite fiber sequences in spectra.
Induction on the remaining directions gives $I_R^{\wedge n}$.
Successive fibers compute the fiber of the initial vertex over the
limit of the punctured cube.  For two directions the relevant identity is
\[
\operatorname{fib}(A\longrightarrow B\times_D C)
\simeq
\operatorname{fib}\bigl(
 \operatorname{fib}(A\to B)\longrightarrow
 \operatorname{fib}(C\to D)\bigr),
\]
obtained by expressing each fiber as a pullback over zero.  Repeated
application gives the assertion for any number of directions, because
finite limits commute.  Thus there is a fiber sequence
\[
I_R^{\wedge n}\xrightarrow{\varphi_{R,n}}\mathbb S
 \longrightarrow
\operatorname*{holim}_{S\in\mathcal P_n^\times}R^{\wedge|S|}.
\]
The chosen nullhomotopy in \Cref{eq:unit-cube-null-cone} gives the
required lift through its suspended fiber.
\end{proof}

\subsection{The quantitative nilpotence input}

The preceding proposition gives Adams filtration at least the number
of restricted factors.  To turn this lift into vanishing,
we use a quantitative form of nilpotence for $MU$.  Its cutoff is
sublinear in the degree and uniform over connective spectra at a fixed
prime.  We state it as a null action on a truncated sphere, the form
needed below for finite sources which may have negative-degree cells.
We write $\tau_{<k}X$ for the Postnikov truncation of $X$ retaining
homotopy groups in degrees below $k$.

\Needspace{6\baselineskip}
\begin{theorem}[Quantitative nilpotence]
\label{thm:quantitative-nilpotence}
Fix a prime $p$, and let $R$ be either $MU_{(p)}$ or the
Brown--Peterson spectrum $BP$ at that prime.
For $k\geq1$, let $f_R(k)$ be the least integer $n\geq0$ for which
\begin{equation}
\label{eq:nilpotence-truncated-action}
I_R^{\wedge n}\wedge\tau_{<k}\mathbb S_{(p)}
 \xrightarrow{\varphi_{R,n}\wedge\operatorname{id}}
\tau_{<k}\mathbb S_{(p)}
\end{equation}
is nullhomotopic.  These integers exist and satisfy
\[
\lim_{k\to\infty}\frac{f_R(k)}{k}=0.
\]
Equivalently, $f_R(k)$ is the least $n$ such that, for every connective
$p$-local spectrum $X$ and every $i<k$, a class in $\pi_iX$ of
$R$-Adams filtration at least $n$ is zero.
\end{theorem}

For $MU_{(p)}$, this is
\cite[Definitions~2.28 and~3.26, Proposition~3.28, and
Theorem~3.30]{Mathew2018}.
For $BP$, the same formulation and its sublinear bound are recorded in
\cite[Definition~B.3, the discussion following Remark~B.4, and
Proposition~B.15]{BurklundHahnSenger2023}.
Both rings are connective, have degree-zero homotopy
$\mathbb Z_{(p)}$, and have finitely generated homotopy in every degree,
as required in those statements.

The truncated action gives the formulation for all connective spectra
directly.
If \Cref{eq:nilpotence-truncated-action} is null, smash it with a
connective $p$-local spectrum $X$.  The map
\[
X\longrightarrow X\wedge\tau_{<k}\mathbb S_{(p)}
\]
is an isomorphism on homotopy in degrees below $k$: its fiber is
$X\wedge\tau_{\geq k}\mathbb S_{(p)}$, which is $k$-connective.
A class with the asserted lift maps to zero on the right and therefore
was zero in $\pi_iX$.

For our application, the key consequence is a null map into the
truncated sphere itself.  Precomposing
\Cref{eq:nilpotence-truncated-action} with the truncation unit
$\eta_k:\mathbb S_{(p)}\to\tau_{<k}\mathbb S_{(p)}$ in its second
factor gives
\begin{equation}
\label{eq:nilpotence-truncated-composite}
\eta_k\circ\varphi_{R,n}=0
\qquad(n\geq f_R(k)).
\end{equation}
The assertion for $n>f_R(k)$ follows by
applying the fiber inclusions to the additional smash factors.

The next observation explains how truncation detects maps from the
finite complexes occurring here.  It uses an upper bound on cell
degrees, without requiring the source to be connective.

\begin{lemma}[Detection by truncation]
\label{lem:cell-truncation}
Suppose that a finite spectrum $K$ has a finite cell presentation with
all cell degrees at most $D$, either in spectra or in $p$-local
spectra.  If $Y$ is $(D+1)$-connective, then
$[K,Y]=0$.  In particular, for every integer $k>D$, the map
\[
[K,\mathbb S_{(p)}]\longrightarrow
[K,\tau_{<k}\mathbb S_{(p)}]
\]
is injective.
\end{lemma}

\begin{proof}
For a single cell in degree $s\leq D$, the mapping group is
$[S^s,Y]=\pi_sY=0$.  Attaching finitely many such cells proves the
first assertion by induction: the cofiber of an attaching step is a
finite wedge of these spheres, and the exact sequence obtained by
mapping into $Y$ makes the mapping group of the enlarged complex zero
when those of the preceding complex and the cofiber are zero.
For the second assertion, apply the first to
$Y=\tau_{\geq k}\mathbb S_{(p)}$ and use the exact sequence obtained
from the fiber sequence
\[
\tau_{\geq k}\mathbb S_{(p)}\longrightarrow\mathbb S_{(p)}
 \longrightarrow\tau_{<k}\mathbb S_{(p)}.
\]
\end{proof}

\subsection{From local vanishing to integral vanishing}

\Cref{prop:compatible-adams-lift} supplies a single lift of filtration
proportional to the number of factors.  We now show that the restricted
product class vanishes once the product is large.  At large primes a
uniform connectivity estimate suffices, leaving only finitely many
prime-dependent nilpotence thresholds.

\Needspace{6\baselineskip}
\begin{lemma}[The unit fiber at large primes]
\label{lem:bp-unit-connectivity}
For every odd prime $p$, the spectrum $BP$ is a retract of $MU_{(p)}$
by unit-preserving maps of spectra, and
\[
I_{BP}\in\mathrm{Sp}_{\geq2p-3}.
\]
Consequently a lift through $(I_{MU})_{(p)}^{\wedge n}$ induces a
lift through a spectrum which is $n(2p-3)$-connective.
\end{lemma}

\begin{proof}
The unit-preserving retract follows, for example, from
the $\mathbb E_4$ splitting of $MU_{(p)}$ in
\cite[Theorem~1.1, the following paragraph, and Section~5]{BasterraMandell2013}.
We use the underlying retract of spectra for ordinary homology and
the forward map under $\mathbb S_{(p)}$ to transfer unit fibers.

Let $C$ be the cofiber of $\mathbb S_{(p)}\to BP$.  The mod-$p$
homology calculation in \cite[Section~5]{BasterraMandell2013} gives
\[
H_0(BP;\mathbb F_p)=\mathbb F_p,\qquad
H_i(BP;\mathbb F_p)=0\quad(0<i<2p-2).
\]
The unit induces the degree-zero isomorphism, and $BP$ is connective.
Its unit cofiber therefore has zero mod-$p$ homology in every degree
below $2p-2$.

The splitting above exhibits $BP$ as a retract of $MU_{(p)}$.
The ordinary integral homology of $MU$ is
$\mathbb Z[b_1,b_2,\ldots]$, with $|b_i|=2i$
\cite[Lecture~7, Proposition~2 and Corollary~3]{Lurie2010}; in particular, it is finitely
generated in every degree.  Taking the retract and then the homology
sequence of the unit cofiber shows that $H_i(C;\mathbb Z_{(p)})$ is
finitely generated in every degree.  For each $i<2p-2$, the universal
coefficient sequence contains an injection
\[
H_i(C;\mathbb Z_{(p)})/pH_i(C;\mathbb Z_{(p)})
 \lhook\joinrel\longrightarrow H_i(C;\mathbb F_p)=0.
\]
Nakayama's lemma therefore gives
$H_i(C;\mathbb Z_{(p)})=0$ for every $i<2p-2$.
The spectrum $C$ is connective.  If it had a nonzero homotopy group
below degree $2p-2$, let $b\geq0$ be the least such degree.
The stable Hurewicz isomorphism in the first nonzero degree would give
$\pi_b(C)\cong H_b(C;\mathbb Z_{(p)})=0$, a contradiction.
Thus $C$ is $(2p-2)$-connective, and
$I_{BP}\simeq\Sigma^{-1}C$ is $(2p-3)$-connective.

The map under the sphere induces a map
$(I_{MU})_{(p)}\to I_{BP}$ commuting with the fiber inclusions.
Smashing $n$ times proves the last assertion, since smash products
add connectivity bounds.
\end{proof}

\begin{proof}[Proof of \Cref{thm:filtered-product}]
Put $e=\dim E$, $n_\ell=\lceil\ell/2\rceil$, and
\[
K_\ell=\Sigma^{-d\ell}\Sigma^\infty U_\ell.
\]
The finite subcomplex $U_\ell\subset E^{\wedge\ell}$ has cell
dimension at most $e\ell$, so $K_\ell$ has upper cell degree
\begin{equation}
\label{eq:filtered-product-cell-bound}
D_\ell=(e-d)\ell.
\end{equation}
The class under consideration is now a map
$\beta_\ell:K_\ell\to\mathbb S$.
By \Cref{prop:compatible-adams-lift}, it lifts through
$I_{MU}^{\wedge n_\ell}$.

If $e<d$, \Cref{lem:cell-truncation} applied directly to the
connective sphere shows that $\beta_\ell=0$ for every $\ell$.
Assume henceforth that $e\geq d$ and put
\[
k_\ell=(e-d)\ell+1.
\]
For a fixed prime $p$, localization of the lift gives a lift through
$I_{MU_{(p)}}^{\wedge n_\ell}$.  If $e>d$,
\Cref{thm:quantitative-nilpotence} gives
\[
f_{MU_{(p)}}(k_\ell)=o(\ell),
\]
so that $n_\ell\geq f_{MU_{(p)}}(k_\ell)$ for all sufficiently
large $\ell$.  If $e=d$, the same inequality eventually holds
because $k_\ell=1$ is fixed and $n_\ell\to\infty$.
By \Cref{eq:nilpotence-truncated-composite}, the composite of
$(\beta_\ell)_{(p)}$ with $\tau_{<k_\ell}\mathbb S_{(p)}$ is
zero.  By \Cref{lem:cell-truncation} and
\Cref{eq:filtered-product-cell-bound}, $(\beta_\ell)_{(p)}=0$.
This proves eventual vanishing for each fixed prime.

To obtain a single threshold, choose an integer $P\geq2$ such that
\begin{equation}
\label{eq:filtered-product-prime-cutoff}
\frac{2p-3}{2}>\max(e-d,0)
\qquad\text{for every prime }p>P.
\end{equation}
For such a prime, \Cref{lem:bp-unit-connectivity} turns the lift of
$(\beta_\ell)_{(p)}$ into a lift through $I_{BP}^{\wedge n_\ell}$.
Its target has connectivity at least
\[
n_\ell(2p-3)\geq\frac\ell2(2p-3)>(e-d)\ell=D_\ell.
\]
The first assertion of \Cref{lem:cell-truncation} makes this lift
zero, for every $\ell$.  Only the finitely many primes at most $P$
remain.  Take the maximum of their eventual thresholds from the
preceding paragraph.  Above that maximum, $\beta_\ell$ vanishes
after localization at every prime.

Because $K_\ell$ is finite, localization of its stable mapping group
agrees with maps from its localization.  One way to see this is the
duality identity
\[
F(K_\ell,\mathbb S)\simeq D K_\ell,
\qquad
(D K_\ell)_{(p)}\simeq
F((K_\ell)_{(p)},\mathbb S_{(p)}),
\]
where $D K_\ell$ denotes the finite Spanier--Whitehead dual.
An element of an abelian group which vanishes in every localization
at a prime is zero.  Indeed, an element of infinite order remains
nonzero in every such localization, and a nonzero element of finite
order remains nonzero at any prime dividing its order.  Applying this
to $[K_\ell,\mathbb S]$ gives $\beta_\ell=0$ integrally, as required.
\end{proof}

\begin{remark}
\label{rem:filtered-product-proportions}
The same proof works with $\lceil\lambda\ell\rceil$ restricted
factors for any fixed $0<\lambda\leq1$.  In
\Cref{prop:compatible-adams-lift}, take
$n=\lceil\lambda\ell\rceil$; in
\Cref{eq:filtered-product-prime-cutoff}, replace $1/2$ by
$\lambda$.  The finite pair affects the eventual value of $\ell$,
but imposes no lower bound on the positive proportion $\lambda$.
\end{remark}
